\documentclass[11pt,a4paper]{article}

\usepackage[utf8]{inputenc}
\usepackage[T1]{fontenc}
\usepackage[english]{babel}

\usepackage{amsmath,amssymb,amsthm,mathtools}
\usepackage[a4paper,margin=2.5cm]{geometry}
\usepackage{lmodern}
\usepackage{microtype}
\usepackage{hyperref}
\usepackage{enumitem}

\title{A Local Structure Theory of the EABC Grid Modulo \(12\)}
\author{Collaborative Draft}
\date{\today}

\newtheorem{theorem}{Theorem}
\newtheorem{lemma}[theorem]{Lemma}
\newtheorem{corollary}[theorem]{Corollary}
\theoremstyle{definition}
\newtheorem{definition}[theorem]{Definition}
\newtheorem{example}[theorem]{Example}
\theoremstyle{remark}
\newtheorem{remark}[theorem]{Remark}

\newcommand{\N}{\mathbb{N}}
\newcommand{\PP}{\mathbb{P}}
\newcommand{\ind}{\mathbf{1}}

\begin{document}
	\maketitle
	
	\begin{abstract}
		We introduce the modulo-\(12\) EABC grid
		\[
		Q_n=(12n-11,\;12n-7,\;12n-5,\;12n-1),
		\qquad n\in\N,
		\]
		and show that it provides a complete local description of small prime constellations above \(3\).
		In this framework, all prime twins \(>3\) lie in exactly two linear channels, all dense prime triplets
		are forced by the singular anchor \(3\), block-crossing prime quadruplets are precisely the double
		activation points of both twin channels, and dense prime quintuplets are exactly the oriented
		extensions of such double points.
	\end{abstract}
	
	\section{The EABC grid}
	
	\begin{definition}[EABC families]
		For \(n\in\N\), define
		\[
		E_n:=12n-11,\qquad
		A_n:=12n-7,\qquad
		B_n:=12n-5,\qquad
		C_n:=12n-1.
		\]
		We call
		\[
		Q_n:=(E_n,A_n,B_n,C_n)
		\]
		the \(n\)-th \emph{EABC block}.
	\end{definition}
	
	\begin{lemma}[Prime decomposition modulo \(12\)]
		Every prime \(p>3\) lies in exactly one of the four families
		\[
		E,\quad A,\quad B,\quad C.
		\]
	\end{lemma}
	
	\begin{proof}
		A prime \(p>3\) is divisible by neither \(2\) nor \(3\). Hence
		\[
		p\equiv 1,5,7,\text{ or }11 \pmod{12}.
		\]
		These four residue classes are parametrized disjointly and exhaustively by
		\[
		12n-11,\quad 12n-7,\quad 12n-5,\quad 12n-1.
		\]
	\end{proof}
	
	\begin{corollary}
		The set of primes \(>3\) decomposes as
		\[
		\PP\setminus\{2,3\}
		=
		(\PP\cap E)\,\dot\cup\,(\PP\cap A)\,\dot\cup\,(\PP\cap B)\,\dot\cup\,(\PP\cap C).
		\]
	\end{corollary}
	
	\section{The two twin channels}
	
	\begin{definition}[Twin channels]
		For each EABC block \(Q_n\), define the internal and external twin channels by
		\[
		T_n^{\mathrm{int}}:=(A_n,B_n)=(12n-7,\;12n-5),
		\]
		\[
		T_n^{\mathrm{ext}}:=(C_n,E_{n+1})=(12n-1,\;12n+1).
		\]
		We further define the \emph{twin signature}
		\[
		\tau(Q_n):=
		\bigl(
		\ind_{T_n^{\mathrm{int}}\text{ is a twin prime pair}},
		\ind_{T_n^{\mathrm{ext}}\text{ is a twin prime pair}}
		\bigr)\in\{0,1\}^2.
		\]
	\end{definition}
	
	\begin{lemma}[Completeness of the twin channels]
		Every twin prime pair \((p,p+2)\) with \(p>3\) lies uniquely in exactly one of the two patterns
		\[
		(12n-7,12n-5)
		\qquad\text{or}\qquad
		(12n-1,12n+1).
		\]
	\end{lemma}
	
	\begin{proof}
		For \(p>3\), a twin prime pair can only have residue pattern
		\[
		(5,7)\quad\text{or}\quad(11,1)\pmod{12}.
		\]
		The first case is exactly
		\[
		(A_n,B_n),
		\]
		and the second case is exactly
		\[
		(C_n,E_{n+1}).
		\]
		Uniqueness follows from the linear parametrization.
	\end{proof}
	
	\begin{corollary}[Classification of twin primes \(>3\)]
		All twin primes \(>3\) are exactly the prime realizations of the two linear forms
		\[
		(12n-7,12n-5)
		\qquad\text{and}\qquad
		(12n-1,12n+1).
		\]
	\end{corollary}
	
	\section{Dense prime triplets}
	
	\begin{definition}[Dense prime triplets]
		A \emph{dense prime triplet} is a prime triple of one of the two forms
		\[
		(p,p+2,p+6)
		\qquad\text{or}\qquad
		(p,p+4,p+6).
		\]
	\end{definition}
	
	\begin{lemma}[Every dense triplet contains \(3\)]
		Every dense prime triplet necessarily contains the prime \(3\).
	\end{lemma}
	
	\begin{proof}
		In either pattern
		\[
		p,\ p+2,\ p+6
		\qquad\text{or}\qquad
		p,\ p+4,\ p+6,
		\]
		at least one of the three numbers is divisible by \(3\). If all three are prime, that number must be \(3\).
	\end{proof}
	
	\begin{corollary}[Classification of dense triplets]
		The only dense prime triplets are
		\[
		(3,5,7)
		\qquad\text{and}\qquad
		(3,7,11).
		\]
	\end{corollary}
	
	\begin{remark}
		Thus dense triplets are not free internal configurations of the EABC grid. They are singular boundary
		phenomena forced by the exceptional prime \(3\).
	\end{remark}
	
	\section{Double points and prime quadruplets}
	
	\begin{definition}[Double point]
		An EABC block \(Q_n\) is called a \emph{double point} if
		\[
		\tau(Q_n)=(1,1).
		\]
		Equivalently, this means that
		\[
		A_n,\ B_n,\ C_n,\ E_{n+1}
		\]
		are all prime.
	\end{definition}
	
	\begin{lemma}[Double points are block-crossing quadruplets]
		An EABC block \(Q_n\) is a double point if and only if
		\[
		(12n-7,\;12n-5,\;12n-1,\;12n+1)
		\]
		is a prime quadruplet.
	\end{lemma}
	
	\begin{proof}
		The condition \(\tau(Q_n)=(1,1)\) means that both
		\[
		(12n-7,12n-5)
		\qquad\text{and}\qquad
		(12n-1,12n+1)
		\]
		are twin prime pairs. This is equivalent to the primality of all four numbers
		\[
		12n-7,\ 12n-5,\ 12n-1,\ 12n+1,
		\]
		which form the pattern
		\[
		p,\ p+2,\ p+6,\ p+8.
		\]
	\end{proof}
	
	\begin{corollary}
		The block-crossing prime quadruplets are exactly the simultaneous activations of the two twin channels.
	\end{corollary}
	
	\section{Dense prime quintuplets}
	
	\begin{definition}[Dense prime quintuplet]
		A \emph{dense prime quintuplet} is a prime constellation of the form
		\[
		(p,p+2,p+6,p+8,p+12).
		\]
	\end{definition}
	
	\begin{lemma}[Two EABC forms of quintuplets]
		Every dense prime quintuplet with \(p>3\) has exactly one of the two forms
		\[
		(A_n,B_n,C_n,E_{n+1},A_{n+1})
		\]
		or
		\[
		(C_n,E_{n+1},A_{n+1},B_{n+1},C_{n+1})
		\]
		for a unique \(n\in\N\).
	\end{lemma}
	
	\begin{proof}
		For \(p>3\), only the starting residues \(5\) and \(11\) modulo \(12\) are possible.
		
		If
		\[
		p\equiv 5 \pmod{12},
		\]
		then \(p=A_n\), hence
		\[
		p+2=B_n,\quad p+6=C_n,\quad p+8=E_{n+1},\quad p+12=A_{n+1}.
		\]
		
		If
		\[
		p\equiv 11 \pmod{12},
		\]
		then \(p=C_n\), hence
		\[
		p+2=E_{n+1},\quad p+6=A_{n+1},\quad p+8=B_{n+1},\quad p+12=C_{n+1}.
		\]
		
		The starting residues \(1\) and \(7\) are impossible, since then one of the subsequent terms would be divisible by \(3\).
	\end{proof}
	
	\begin{lemma}[Every dense quintuplet contains exactly one double point]
		Every dense prime quintuplet contains exactly one block-crossing prime quadruplet.
	\end{lemma}
	
	\begin{proof}
		In the first form
		\[
		(A_n,B_n,C_n,E_{n+1},A_{n+1}),
		\]
		the first four terms form the quadruplet.
		
		In the second form
		\[
		(C_n,E_{n+1},A_{n+1},B_{n+1},C_{n+1}),
		\]
		the first four terms form the quadruplet.
		
		A dense quintuplet has only five entries, so it cannot contain two distinct quadruplets of this type.
	\end{proof}
	
	\begin{corollary}
		Every dense prime quintuplet is an oriented extension of a unique double point.
	\end{corollary}
	
	\section{Main theorem}
	
	\begin{theorem}[Local completeness of the EABC grid]
		The EABC grid
		\[
		Q_n=(12n-11,\;12n-7,\;12n-5,\;12n-1)
		\]
		provides a complete local description of small prime constellations above \(3\):
		
		\begin{enumerate}[label=\textup{(\roman*)}]
			\item Every prime \(p>3\) lies uniquely in one of the four families \(E,A,B,C\).
			\item Every twin prime pair \(>3\) lies uniquely in one of the two channels
			\[
			(A_n,B_n)
			\qquad\text{or}\qquad
			(C_n,E_{n+1}).
			\]
			\item Every dense prime triplet is one of the two singular \(3\)-anchored cases
			\[
			(3,5,7),\qquad (3,7,11).
			\]
			\item Every double point \(\tau(Q_n)=(1,1)\) is exactly a block-crossing prime quadruplet
			\[
			(12n-7,12n-5,12n-1,12n+1).
			\]
			\item Every dense prime quintuplet is exactly an oriented extension of such a double point, hence has one of the two forms
			\[
			(A_n,B_n,C_n,E_{n+1},A_{n+1})
			\qquad\text{or}\qquad
			(C_n,E_{n+1},A_{n+1},B_{n+1},C_{n+1}).
			\]
		\end{enumerate}
		
		Therefore the full local twin-, triplet-, quadruplet-, and quintuplet-structure of the prime grid modulo \(12\) is completely described.
	\end{theorem}
	
	\begin{proof}
		Statements \textup{(i)} and \textup{(ii)} follow from the first two lemmas. Statement \textup{(iii)} is the triplet corollary. Statement \textup{(iv)} is the quadruplet lemma. Statement \textup{(v)} follows from the two quintuplet lemmas.
	\end{proof}
	
	\begin{remark}
		In this picture, the prime quadruplet is the simultaneous activation of both twin channels, while the prime quintuplet is the first genuinely directed continuation of such a double point.
	\end{remark}
	
\end{document}